% =============================================================================
%  lec10.tex —— 第 10 讲 极限与余极限：交换性（底本 PDF 第 80–87 页）
% =============================================================================
\lecture{10}{极限与余极限：交换性}
\begin{center}\small\color{black!55}（底本 PDF 第 80–87 页）\end{center}

上一讲中，对图 $u\colon K_1\times K_2\to\cC$，我们在假定两侧均存在时构造了典范态射
\[
\colim_{y\in K_2}\lim_{K_1}u(-,y)\to\lim_{x\in K_1}\colim_{K_2}u(x,-).
\]
本讲讨论上述态射可逆的若干情形。

\section{有限图}

\begin{definition}\label{L10-1.1}
若单纯集 $K$ 只有有限多个非退化单形，则称 $K$ 为\textbf{有限的}。

若 $\infty$-范畴 $K$ 可由一个作为有限单纯集的拟范畴表示，则称 $K$ 为\textbf{本质有限的}。
\end{definition}

\nitem\label{L10-1.2} 回忆：我们称普通范畴 $K$ 本质有限，当且仅当 $K$ 中对象的等价类构成的集合与 $K$ 中态射构成的集合都有限。

\begin{exercise}\label{L10-1.3}
设 $K$ 是本质有限的 $\infty$-范畴。证明 $hK$ 是本质有限的普通范畴。
\end{exercise}

\begin{exercise}\label{L10-1.4}
设 $K$ 是普通范畴。若 $K$ 作为 $\infty$-范畴是本质有限的，则它作为普通范畴是本质有限的。
\end{exercise}

\begin{warning}\label{L10-1.5}
逆命题不成立：即便是有限的普通范畴，也可能作为 $\infty$-范畴不是本质有限的。
\end{warning}

\begin{exercise}\label{L10-1.6}
设 $\mathrm{Idem}$ 是由下述条件定义的普通范畴：
\begin{itemize}[leftmargin=2.2em,itemsep=0.1em]
\item 有唯一的对象 $\ast$；
\item $\Hom(\ast,\ast)\coloneqq\{\id,e\}$，其中 $e\circ e=e$。
\end{itemize}
证明 $N_{\bullet}(\mathrm{Idem})$ 不等价于一个有限的拟范畴。
\end{exercise}

\begin{exercise}\label{L10-1.7}
设 $\mathbb{B}(\Z/2\Z)$ 是由下述条件定义的普通范畴：
\begin{itemize}[leftmargin=2.2em,itemsep=0.1em]
\item 有唯一的对象 $\ast$；
\item $\Hom(\ast,\ast)\coloneqq\{\id,f\}$，其中 $f\circ f=\id$。
\end{itemize}
证明 $N_{\bullet}(\mathbb{B}(\Z/2\Z))$ 不等价于一个有限的拟范畴。
\end{exercise}

\nitem\label{L10-1.8} 不过，对有限偏序集而言，上述异常现象不会出现。

\begin{exercise}\label{L10-1.9}
设 $J$ 是有限偏序集，视为普通范畴。证明 $N_{\bullet}(J)$ 是有限单纯集。
\end{exercise}

\begin{proposition}\label{L10-1.10}
（Ker.02NB）对任意单纯集 $K$，存在偏序集 $J$ 与初始态射 $N_{\bullet}(J)\to K$。若 $K$ 有限，则可取 $J$ 有限。
\end{proposition}

\nitem\label{L10-1.11}\textbf{推论—定义}\quad 设 $\cC$ 是 $\infty$-范畴。下列条件等价：
\begin{itemize}[leftmargin=2.2em,itemsep=0.1em]
\item $\cC$ 容许以有限单纯集为指标的极限。
\item $\cC$ 容许以本质有限的 $\infty$-范畴为指标的极限。
\item $\cC$ 容许以有限偏序集为指标的极限。
\end{itemize}
若 $\cC$ 满足上述条件，则称 $\cC$ \textbf{容许有限极限}。

设 $F\colon\cC\to\cD$ 是 $\infty$-范畴之间的函子，且 $\cC$ 容许有限极限。下列条件等价：
\begin{itemize}[leftmargin=2.2em,itemsep=0.1em]
\item $F$ 保持以有限单纯集为指标的极限。
\item $F$ 保持以本质有限的 $\infty$-范畴为指标的极限。
\item $F$ 保持以有限偏序集为指标的极限。
\end{itemize}
若 $F$ 满足上述条件，则称 $F$ 是\textbf{左正合的}。

对偶地，设 $F\colon\cC\to\cD$ 是 $\infty$-范畴之间的函子，且 $\cC$ 容许有限余极限。若 $F$ 保持有限余极限，则称 $F$ 是\textbf{右正合的}。

\begin{warning}\label{L10-1.12}
设 $\cC$ 是容许有限极限的 $\infty$-范畴。它未必容许以有限普通范畴为指标的极限。例如，$\cC$ 可能不是幂等完备的。
\end{warning}

\begin{proposition}\label{L10-1.13}
设 $\cC$ 是 $\infty$-范畴。下列条件等价：
\begin{enumerate}[label=(\roman*),leftmargin=2.4em,itemsep=0.1em]
\item $\cC$ 容许有限极限。
\item $\cC$ 容许纤维积与终对象。
\item $\cC$ 容许有限积与等化子。
\end{enumerate}
对左正合函子，我们也有类似的结果。
\end{proposition}

\begin{proof}[证明梗概]
蕴含 (i)$\Rightarrow$(ii)$\Leftrightarrow$(iii) 是标准的。为证明 (ii)$\Rightarrow$(i)，设 $u\colon K\to\cC$ 是有限图。我们需要证明 $\lim u$ 存在。当 $K=\varnothing$ 时，这就是 $\cC$ 的终对象。若 $K$ 非空，我们可在 $\Set_\Delta$ 中找到推挤方块
\begin{center}
\begin{tikzcd}[column sep=3.2em,row sep=2.8em]
\partial\Delta^n \arrow[r] \arrow[d,"{\subset}"{sloped=false}] & K' \arrow[d,"{\subset}"{sloped=false}] \\
\Delta^n \arrow[r,"{\subset}"{sloped=false}] & K
\end{tikzcd}
\end{center}
注意这同时也是 $\Set^{\mathrm{Joyal}}_\Delta$ 中的同伦推挤方块。由图的分解定理（Ker.03DB），我们有
\[
\lim u\xrightarrow{\ \simeq\ }\lim u|_{K'}\times_{\lim u|_{\partial\Delta^n}}\lim u|_{\Delta^n},
\]
其中若靶存在，则源存在。由归纳，我们归约为证明 $K=\Delta^n$ 时的断言，而这由 $\Delta^0\to\Delta^n$ 是初始态射这一事实得到。
\end{proof}

\nitem\label{L10-1.14} 事实上，对上述论证加以细化便可得：

\begin{proposition}\label{L10-1.15}
（HTT.4.4.2.6, 4.4.3.2）设 $\cC$ 是 $\infty$-范畴。下列条件等价：
\begin{enumerate}[label=(\roman*),leftmargin=2.4em,itemsep=0.1em]
\item $\cC$ 容许小极限。
\item $\cC$ 容许纤维积与小积。
\item $\cC$ 容许小积与等化子。
\end{enumerate}
对保持小极限的函子，我们也有类似的结果。
\end{proposition}

\section{滤过的 $\infty$-范畴}

\begin{definition}\label{L10-2.1}
若任意有限图 $u\colon K\to\cC$ 都能扩张为图 $K^{\triangleright}\to\cC$，则称 $\infty$-范畴 $\cC$ 是\textbf{滤过的}（filtered）。

对偶地，若任意有限图 $u\colon K\to\cC$ 都能扩张为图 $K^{\triangleleft}\to\cC$，则称 $\cC$ 是\textbf{余滤过的}。
\end{definition}

\nitem\label{L10-2.2} 注意 $\cC$ 滤过当且仅当 $\cC^{\op}$ 余滤过。因此下文集中讨论滤过的 $\infty$-范畴。

\begin{exercise}\label{L10-2.3}
滤过性在等价下不变。
\end{exercise}

\nitem\label{L10-2.4} 一个归纳论证（它与命题~\ref{L10-1.13} 的证明类似，但更为精细）给出下述结果：

\begin{lemma}\label{L10-2.5}
（Ker.02Q0）$\infty$-范畴 $\cC$ 滤过，当且仅当任意图 $u\colon\partial\Delta^n\to\cC$（$n\ge 0$\footnote{按定义，$\partial\Delta^0=\varnothing$。}）都能扩张为图 $(\partial\Delta^n)^{\triangleright}\to\cC$。
\end{lemma}

\begin{remark}\label{L10-2.6}
注意 $(\partial\Delta^n)^{\triangleright}\simeq\Lambda^{n+1}_{n+1}$。
\end{remark}

\nitem\label{L10-2.7} 回忆：我们称普通范畴 $\cC$ 滤过，若
\begin{enumerate}[label=(\arabic*),leftmargin=2.4em,itemsep=0.1em,start=0]
\item 它非空。
\item 对任意对象 $x_1,x_2\in\cC$，存在带态射 $x_i\to x$ 的对象 $x$。
\item 对 $\cC$ 中任意图 $x\rightrightarrows y$，存在 $y\to z$，使得两个复合 $x\rightrightarrows y\to z$ 相等。
\end{enumerate}
注意条件 $(i)$ 对应 $K=\partial\Delta^i$ 时的扩张问题。

\begin{proposition}\label{L10-2.8}
（Ker.02PS）$\infty$-范畴 $\cC$ 滤过，当且仅当它满足下述条件：
\begin{enumerate}[label=(\arabic*),leftmargin=2.4em,itemsep=0.1em,start=0]
\item 它非空。
\item 对任意对象 $x_1,x_2\in\cC$，存在带态射 $x_i\to x$ 的对象 $x$。
\item 对 $\cC$ 中任意对象 $x,y$ 与任意图 $\partial\Delta^n\to\operatorname{Maps}(x,y)$（$n\ge 1$），存在态射 $y\to z$，使得复合
\[
\partial\Delta^n\to\operatorname{Maps}(x,y)\to\operatorname{Maps}(x,z)
\]
零伦，即等价于常值图\footnote{练习：对 Kan 复形 $K$，态射 $\partial\Delta^n\to K$ 零伦当且仅当它通过 $\Delta^n$ 分解。}。
\end{enumerate}
\end{proposition}

\begin{remark}\label{L10-2.9}
粗略地说，对 $n\ge 2$，扩张问题
\begin{center}
\begin{tikzcd}[column sep=3.6em,row sep=3.0em]
\partial\Delta^n \arrow[r,"u"] \arrow[d] & \cC \\
\Delta^n \arrow[ur,dotted,-{Latex}] &
\end{tikzcd}
\end{center}
对应于扩张问题
\begin{center}
\begin{tikzcd}[column sep=4.8em,row sep=3.0em]
\partial\Delta^{n-1} \arrow[r] \arrow[d] & \operatorname{Maps}(x,y) \\
\Delta^{n-1} \arrow[ur,dotted,-{Latex}] &
\end{tikzcd}
\end{center}
其中 $x=u(0)$、$y=u(n)$。
\end{remark}

\begin{exercise}\label{L10-2.10}
设 $\cC$ 是滤过的 $\infty$-范畴。证明 $h\cC$ 是滤过的普通范畴。
\end{exercise}

\begin{exercise}\label{L10-2.11}
设 $\cC$ 是普通范畴。证明：$\cC$ 作为普通范畴滤过，当且仅当它作为 $\infty$-范畴滤过。
\end{exercise}

\begin{example}\label{L10-2.12}
若 $\infty$-范畴 $\cC$ 含有终对象，则它滤过。
\end{example}

\begin{example}\label{L10-2.13}
若 $\infty$-范畴 $\cC$ 容许有限余极限，则它滤过。
\end{example}

\begin{exercise}\label{L10-2.14}
范畴 $\mathrm{Idem}$ 滤过。
\end{exercise}

\section{滤过余极限与有限极限}

\begin{theorem}\label{L10-3.1}
（Ker.05XS）设 $\cC$ 是小 $\infty$-范畴。则 $\cC$ 滤过，当且仅当函子
\[
\colim\colon\Fun(\cC,\Grpd_\infty)\to\Grpd_\infty
\]
保持有限极限。
\end{theorem}

\begin{warning}\label{L10-3.2}
若把 $\Grpd_\infty$ 换成一般的 $\infty$-范畴，该结论可能不成立。
\end{warning}

\begin{exercise}\label{L10-3.3}
考虑离散拓扑空间 $\Z$ 及其单点紧化 $\Z\cup\{\infty\}$\footnote{$\Z\cup\{\infty\}$ 的闭子集要么有限，要么包含 $\infty$。}。考虑由 $\Z\cup\{\infty\}$ 的闭子集构成的偏序集 $\cP$。证明 $\cP$ 中的滤过余极限未必与有限极限交换。
\end{exercise}

\begin{corollary}\label{L10-3.4}
设 $F\colon\cC\to\cD$ 是 $\infty$-范畴之间的终结函子。若 $\cC$ 滤过，则 $\cD$ 也滤过。
\end{corollary}

\begin{challenge}\label{L10-3.5}
你能给出上述推论的一个证明而不使用定理~\ref{L10-3.1} 吗？
\end{challenge}

\nitem\label{L10-3.6} 回忆：偏序集 $J$ 称为\textbf{有向的}，若它的任意有限子集都在 $J$ 中有上界。

\begin{proposition}\label{L10-3.7}
（Ker.02QA）设 $\cC$ 是 $\infty$-范畴。下列条件等价：
\begin{itemize}[leftmargin=2.2em,itemsep=0.1em]
\item $\cC$ 滤过。
\item 存在有向偏序集 $J$ 与终结函子 $J\to\cC$。
\end{itemize}
\end{proposition}

\begin{proposition}\label{L10-3.8}
设 $\cC$ 是滤过的 $\infty$-范畴，$u\colon K\to\cC$ 是任意有限图。则遗忘函子 $\cC_{u/}\to\cC$ 终结。
\end{proposition}

\begin{proof}
由 Quillen 定理 A，我们只需证明对任意 $x\in\cC$，$\cC_{u/}\times_\cC\cC_{x/}$ 弱可缩。我们有
\[
\cC_{u/}\times_\cC\cC_{x/}\simeq\cC_{u\sqcup x/},
\]
其中 $u\sqcup x\colon K\sqcup\Delta^0\to\cC$ 是 $u$ 与 $x$ 的不交并。现在断言由下面两个练习推出。
\end{proof}

\begin{exercise}\label{L10-3.9}
滤过的 $\infty$-范畴弱可缩。提示：$K\times\Delta^1\to K^{\triangleright}$。
\end{exercise}

\begin{proposition}\label{L10-3.10}
设 $\cC$ 是 $\infty$-范畴。下列条件等价：
\begin{enumerate}[label=(\roman*),leftmargin=2.4em,itemsep=0.1em]
\item $\cC$ 滤过。
\item 对任意有限图 $u\colon K\to\cC$，$\infty$-范畴 $\cC_{u/}$ 滤过。
\item 对任意有限图 $u\colon K\to\cC$，$\infty$-范畴 $\cC_{u/}$ 弱可缩。
\item 对任意有限 $K\in\Set_\Delta$，对角函子 $\cC\to\Fun(K,\cC)$，$x\mapsto\underline{x}$ 终结。
\end{enumerate}
\end{proposition}

\begin{proof}[证明梗概]
蕴含 (i)$\Rightarrow$(ii)$\Rightarrow$(iii)$\Rightarrow$(i) 留作练习。我们将证明 (ii)$\Leftrightarrow$(iii)。由 Quillen 定理 A，(ii) 等价于：对任意 $u\colon K\to\cC$，纤维积
\[
\cC\times_{\Fun(K,\cC)}\Fun(K,\cC)^{u/}
\]
弱可缩。容易看出上述纤维积等价于 $\cC^{u/}$，这就推出断言。
\end{proof}

\section{筛图}

\begin{definition}\label{L10-4.1}
设 $K$ 是单纯集。若对任意有限集 $I$，对角态射 $K\to K^I$ 终结，则称 $K$ 是\textbf{筛的}（sifted）。

对偶地，若对任意有限集 $I$，对角态射 $K\to K^I$ 初始，则称 $K$ 是\textbf{余筛的}。
\end{definition}

\yizhu{术语 sifted 本义为“筛过的”，此处译作“筛的”，其后简称“筛”；相应地 sifted colimit 译作“筛余极限”，与 filtered（滤过的）相区分。\textbf{[来源：B]}}

\nitem\label{L10-4.2} 注意 $\cC$ 筛当且仅当 $\cC^{\op}$ 余筛。因此下文集中讨论筛的 $\infty$-范畴。

\begin{exercise}\label{L10-4.3}
滤过性在等价下不变。
\end{exercise}

\yizhu{练习 4.3 的原文为“Being filtered is invariant under equivalences”，与练习 2.3 逐字重复；按本节主题（筛性）以及它与练习 2.3 的对应关系，此处原文当作“Being sifted”。译文照录原文。\textbf{[来源：I]}}

\begin{example}\label{L10-4.4}
任意滤过的 $\infty$-范畴都是筛的。
\end{example}

\begin{exercise}\label{L10-4.5}
证明单纯集 $K$ 筛，当且仅当它非空且 $K\to K\times K$ 终结。提示：若 $K\to K\times K$ 是弱同伦等价且 $K\neq\varnothing$，则 $K$ 弱可缩。
\end{exercise}

\begin{corollary}\label{L10-4.6}
$\infty$-范畴 $\cC$ 筛，当且仅当它非空，且对任意一对对象 $x,y\in\cC$，$\cC_{x/}\times_\cC\cC_{y/}$ 弱可缩。
\end{corollary}

\begin{warning}\label{L10-4.7}
在经典范畴论中，普通范畴 $\cC$ 称为筛的，若它非空，且对任意一对对象 $x,y\in\cC$，$\cC_{x/}\times_\cC\cC_{y/}$ 连通。筛的普通范畴作为 $\infty$-范畴可能不筛。
\end{warning}

\begin{exercise}\label{L10-4.8}
设 $\Delta_{\le 1}\subset\Delta$ 是由 $[0]$ 与 $[1]$ 构成的全子范畴。证明 $\Delta_{\le 1}^{\op}$ 作为普通范畴筛，但作为 $\infty$-范畴不筛。
\end{exercise}

\begin{exercise}\label{L10-4.9}
证明 $\Delta^{\op}$ 作为 $\infty$-范畴筛。
\end{exercise}

\begin{definition}\label{L10-4.10}
$\Delta^{\op}$-指标图的余极限称为该图的\textbf{几何实现}。$\Delta$-指标图的极限称为该图的\textbf{全化}。
\end{definition}

\nitem\label{L10-4.11} 众所周知，几何实现与滤过余极限几乎生成全部的筛图。见 §A.10。

\begin{proposition}\label{L10-4.12}
设 $F\colon\cC\to\cD$ 是 $\infty$-范畴之间的终结函子。若 $\cC$ 筛，则 $\cD$ 也筛。
\end{proposition}

\begin{proof}
对任意有限集 $I$，考虑交换图
\begin{center}
\begin{tikzcd}[column sep=3.8em,row sep=3.2em]
\cC \arrow[r,"F"] \arrow[d,"\delta_1"'] & \cD \arrow[d,"\delta_2"] \\
\cC^I \arrow[r,"F^I"'] & \cD^I
\end{tikzcd}
\end{center}
由于 $F$ 终结，$F^I$ 也终结（第 7 讲命题 3.16）。由于 $\cC$ 筛，$\delta_1$ 终结。于是 $\delta_2$ 也终结，如所愿望（第 7 讲练习 3.13）。
\end{proof}

\begin{exercise}\label{L10-4.13}
设 $K$ 是筛的单纯集。设 $\cC$、$\cD$、$\cE$ 是容许 $K$-指标余极限的 $\infty$-范畴，$F\colon\cC\times\cD\to\cE$ 是在每个变量上都保持 $K$-指标余极限的函子\footnote{即对任意 $c\in\cC$ 与 $d\in\cD$，$F(c,-)$ 与 $F(-,d)$ 都保持 $K$-指标余极限。}。证明 $F$ 保持 $K$-指标余极限。
\end{exercise}

\section{筛余极限与有限积}

\begin{theorem}\label{L10-5.1}
（HTT.5.5.8.11, Ker.05XM）设 $\cC$ 是小 $\infty$-范畴。则 $\cC$ 筛，当且仅当函子
\[
\colim\colon\Fun(\cC,\Grpd_\infty)\to\Grpd_\infty
\]
保持有限积。
\end{theorem}

\section{余极限的泛性}

\begin{construction}\label{L10-6.1}
设 $\cC$ 是容许纤维积的 $\infty$-范畴。设 $X\to Y$ 是 $\Grpd_\infty$ 中的函子。与经典范畴论一样，我们有伴随
\begin{equation}\label{L10-6.1eq}
(\Grpd_\infty)_{/X}\rightleftarrows(\Grpd_\infty)_{/Y},
\end{equation}
其中：
\begin{itemize}[leftmargin=2.2em,itemsep=0.1em]
\item 左伴随与遗忘函子相容；
\item 右伴随由 $-\times_Y X$ 给出。
\end{itemize}
\end{construction}

\begin{remark}\label{L10-6.2}
或者，上述右伴随可通过单纯模型范畴 $\Set^{\mathrm{KQ}}_\Delta$ 构造。
\end{remark}

\begin{theorem}\label{L10-6.3}
（Ker.05V5）函子
\[
-\times_Y X\colon(\Grpd_\infty)_{/Y}\to(\Grpd_\infty)_{/X}
\]
保持小余极限。
\end{theorem}

\begin{exercise}\label{L10-6.4}
设 $u\colon K\to\Grpd_\infty$，$i\mapsto X_i$ 是小图，记 $X\coloneqq\colim u$。证明存在典范同构
\[
\colim_{i\in K}(X_i\times_X Y)\simeq Y.
\]
\end{exercise}

\appendixsection{滤过余极限}

\begin{proposition}\label{L10-A.1}
设 $\cC$ 是小滤过普通范畴。则任意余极限图 $\bar u\colon\cC^{\triangleright}\to\Set_\Delta$ 也是 $\Set^{\mathrm{Joyal}}_\Delta$ 中的同伦余极限图。
\end{proposition}

\nitem\label{L10-A.2} 我们介绍命题~\ref{L10-A.1} 的两个应用。第一个应用是 $\infty$-范畴的滤过余极限的一个描述。

\begin{exercise}\label{L10-A.3}
函子 $\mathrm{QCat}\to\Cat_\infty$ 保持小的滤过 $\infty$-余极限。
\end{exercise}

\begin{exercise}\label{L10-A.4}
设 $K$ 是小滤过 $\infty$-范畴，$K\to\Cat_\infty$，$i\mapsto\cC_i$ 是图。证明 $\cC\coloneqq\colim_K\cC_i$ 中的任意对象都等价于某个 $i\in K$、$x\in\cC_i$ 的 $\mathrm{ins}_i(x)$。
\end{exercise}

\begin{exercise}\label{L10-A.5}
设 $K$ 是小滤过 $\infty$-范畴，$K\to\Cat_\infty$，$i\mapsto\cC_i$ 是图。对 $i,j\in K$ 与 $x\in\cC_i$、$y\in\cC_j$，证明
\[
\operatorname{Maps}_\cC(\mathrm{ins}_i(x),\mathrm{ins}_j(y))\xleftarrow{\ \simeq\ }\colim_{k\in K_{(i,j)/}}\operatorname{Maps}_{\cC_k}(x_k,y_k),
\]
其中
\begin{itemize}[leftmargin=2.2em,itemsep=0.1em]
\item $K_{(i,j)/}$ 是关于 $\Delta^0\sqcup\Delta^0\xrightarrow{(i,j)}K$ 的余切片 $\infty$-范畴；
\item 对 $k\in K_{(i,j)/}$，$x_k$ 是 $x$ 在函子 $\cC_i\to\cC_k$ 下的像，$y_k$ 类似。
\end{itemize}
\end{exercise}

\nitem\label{L10-A.6} 第二个应用是把一般的余极限分解为有限余极限的滤过余极限。

\begin{exercise}\label{L10-A.7}
设 $K$ 是任意单纯集。证明 $K$ 的有限单纯子集构成的偏序集 $\mathrm{Fin}(K)$ 滤过，且
\[
K\simeq\hocolim_{J\in\mathrm{Fin}(K)}J.
\]
\end{exercise}

\begin{exercise}\label{L10-A.8}
对任意图 $u\colon K\to\cC$，有典范等价
\[
\colim_{J\in\mathrm{Fin}(K)}\colim_J u|_J\xrightarrow{\ \simeq\ }\colim_K u,
\]
其中右端在左端存在时存在。
\end{exercise}

\begin{corollary}\label{L10-A.9}
$\infty$-范畴 $\cC$ 容许小余极限，当且仅当它容许小的滤过余极限与有限余极限。对这样的 $\cC$，函子 $F\colon\cC\to\cD$ 保持小余极限，当且仅当它保持滤过余极限与有限余极限。
\end{corollary}

\nitem\label{L10-A.10} 上述结果说明
\[
\text{small colimits}=\text{small filtered colimits}+\text{finite colimits}.
\]
事实上，我们还有：
\[
\text{small colimits}=\text{small sifted colimits}+\text{finite coproducts}.
\]
\[
\text{small colimits}=\text{small filtered colimits}+\text{geometric realization}+\text{finite coproducts}.
\]
见 HA.1.3.3.10。

\nitem\label{L10-A.11} \textbf{建议阅读}\quad Ker.03DD.

\appendixsection{余极限的强泛性}

\begin{theorem}\label{L10-B.1}
设 $K\to\Grpd_\infty$，$i\mapsto X_i$ 是小图。考虑 $X\coloneqq\colim_{i\in K}X_i$。则存在典范等价
\begin{equation}\label{L10-B.1eq}
(\Grpd_\infty)_{/X}\to\lim_{i\in K}(\Grpd_\infty)_{/X_i},
\end{equation}
其中每个求值函子是
\[
-\times_X X_i\colon(\Grpd_\infty)_{/X}\to(\Grpd_\infty)_{/X_i}.
\]
\end{theorem}

\begin{exercise}\label{L10-B.2}
利用定理~\ref{L10-6.3} 证明 \eqref{L10-B.1eq} 全忠实。
\end{exercise}

\begin{exercise}\label{L10-B.3}
证明：若把 $\Grpd_\infty$ 换成 $\Set$，定理~\ref{L10-6.3} 仍然成立，但定理~\ref{L10-B.1} 将不成立。
\end{exercise}

\nitem\label{L10-B.4} \textbf{建议阅读}\quad Ker.05SB.
